% Worked solutions for Listas/lista-r2c-financial-planning-forecasting.md (15 exercises).
% Compile: pdflatex twice; check: pdffonts (Type 1 only, uni = yes).
% Every number is asserted in lista-r2c-financial-planning-forecasting-solutions_check.py.
\documentclass[11pt,a4paper]{article}
\usepackage[utf8]{inputenc}
\usepackage[T1]{fontenc}
\usepackage{lmodern}   % Type 1 vector fonts (replaces PK bitmaps)
\usepackage{cmap}      % ToUnicode CMap -> copyable text and accents
\usepackage[brazilian,english]{babel}
\usepackage{amsmath,amssymb,amsfonts}
\usepackage[a4paper,margin=2.2cm]{geometry}
\usepackage{booktabs}
\usepackage{enumitem}
\usepackage{fancyhdr}
\usepackage{xcolor}
\usepackage{microtype}
\usepackage{tikz}

\pagestyle{fancy}\fancyhf{}
\lhead{\footnotesize Finan\c{c}as II --- Corporativas (MPE Insper)}
\rhead{\footnotesize List R2c (financial planning) --- Worked solutions}
\cfoot{\footnotesize\thepage}
\renewcommand{\headrulewidth}{0.3pt}
\setlength{\headheight}{14pt}
\setlist{topsep=2pt,itemsep=1pt}

\newcommand{\exsec}[1]{\section*{\color{blue!55!black}#1}\addcontentsline{toc}{section}{#1}}
\newcommand{\notation}{\par\noindent\textbf{Notation.}}
\newcommand{\step}[2]{\par\medskip\noindent\textbf{Step #1 --- #2.}~}
\newcommand{\intuition}{\par\medskip\noindent\textbf{Intuition.}~}
\newcommand{\bookans}{\par\medskip\noindent\textbf{Book answer.}~}
\newcommand{\refbib}{\par\smallskip\noindent\textbf{Ref:}~}
\newcommand{\asked}{\par\medskip\noindent\textbf{What is asked.}~}
\newcommand{\approach}{\par\smallskip\noindent\textbf{Approach.}~}
\newcommand{\answer}[1]{\par\smallskip\noindent\fbox{\textbf{Answer #1.}}~}
\newcommand{\pitfalls}{\par\medskip\noindent\textbf{Common mistakes.}}
\newcommand{\EFN}{\mathit{EFN}}

\usepackage[hyphens]{url}
\begin{document}
\begin{center}{\Large\bfseries Exercise list R2c --- Financial planning and forecasting for valuation}\\[4pt]{\large Worked solutions}\end{center}

\noindent Statements are in \url{Listas/lista-r2c-financial-planning-forecasting.md}; each exercise is
cited by number and printed page and is not restated. Short answers:
\url{lista-r2c-financial-planning-forecasting-solutions-short.pdf}. Editions are the local copies:
Berk = Berk, DeMarzo \& Harford, \emph{Fundamentals of Corporate Finance}, 2nd ed.\ (2011), Ch.~18;
Higgins = Koski, Higgins \& Mitton, \emph{Analysis for Financial Management}, 13th ed., Ch.~3--4;
Ross = Ross, Westerfield, Jordan \& Lamb, \emph{Fundamentos de Administra\c{c}\~ao Financeira}, 9th ed.,
Cap.~4 (``NAF'' = \emph{necessidade de aporte financeiro} = EFN). Amounts carry the units of the
statement (dollars, \$ thousands or \$ millions). Numbers marked ``illustration'' are not part of the
statement.

\tableofcontents

% =============================================================================
\exsec{1. Cost of goods sold by the percent of sales method --- Berk 18.1, p.~558}

\asked Next year's cost of goods sold (COGS) when sales rise from \$100,000 to \$110,000, assuming COGS
keeps this year's ratio to sales.

\approach The percent of sales method (Berk \S18.2) assumes that an item that moves with activity keeps
the same fraction of sales. So we measure the fraction today and apply it to forecast sales. This is
the first line of any pro forma income statement.

\notation
\begin{itemize}
  \item $S_0=100{,}000$, $S_1=110{,}000$: sales this year and next year (\$).
  \item $C_0=72{,}000$: COGS this year (\$); $c=C_0/S_0$: COGS as a fraction of sales; $C_1$: forecast COGS.
\end{itemize}

\step{1}{Measure the ratio} We divide this year's COGS by this year's sales:
\begin{align}
  c &= \frac{C_0}{S_0} = \frac{72{,}000}{100{,}000} = 0.72. \label{r2ce1-c}
\end{align}

\step{2}{Apply it to forecast sales} The ratio is assumed constant, so
\begin{align}
  C_1 &= c\,S_1 = (0.72)(110{,}000) = 79{,}200. \label{r2ce1-C1}
\end{align}
\answer{} Forecast COGS $=\$79{,}200$ (72\% of \$110,000).

\intuition Holding COGS at 72\% of sales assumes the gross margin (28\%) is stable: each extra dollar of
sales brings 72 cents of extra cost. It is a good first guess for a mature business; a firm with
fixed production costs would see COGS grow more slowly than sales.
\pitfalls
\begin{itemize}
  \item Growing COGS by a different rate than sales (e.g.\ by last year's COGS growth) --- the method ties
        every proportional item to the \emph{sales} forecast.
  \item Applying the ratio to this year's sales instead of next year's.
\end{itemize}
\refbib{Berk \S18.2, ``Percent of Sales Method'', pp.~537--538; Example 18.1, p.~539.}

% =============================================================================
\exsec{2. Reading external financing required --- Higgins 3.1 and 3.4, p.~103}

\asked (3.1) What a \emph{negative} external financing required means. (3.4a) Whether a larger
receivables forecast raises external funding required. (3.4b) Whether cash flow forecasts usually show
larger needs than pro forma statements. (3.4c) Whether an annual forecast with no external need rules
out cash shortfalls during the year.

\approach In a pro forma balance sheet, external funding required (EFN) is the \emph{plug}: the amount
that makes assets equal liabilities plus equity. Every question follows from the identity that defines
it and from the date at which the identity is evaluated.

\notation
\begin{itemize}
  \item $TA$: projected total assets; $L$: projected liabilities (without new financing); $E$: projected equity.
  \item $\EFN = TA-(L+E)$: external funding required at the forecast date.
\end{itemize}

\step{1}{The defining identity} By construction,
\begin{align}
  \EFN &= TA-(L+E). \label{r2ce2-efn}
\end{align}
A positive value means projected uses (assets) exceed projected sources: money must be raised.

\step{2}{(3.1) Negative EFN} If $\EFN<0$, sources exceed uses: the company will generate more cash than
its plan absorbs. Illustration: projected assets $1{,}000$ with cash at its minimum of $100$, and
$L+E=1{,}040$; then $\EFN=1{,}000-1{,}040=-40$. Setting the plug to zero and adding the $40$ to cash gives
cash $=100+40=140$ and assets $=900+140=1{,}040=L+E$: the balance sheet balances.
\answer{(3.1)} A negative EFN is a projected cash \emph{surplus} above the desired minimum: cash (or
debt repayment, dividends, buybacks) absorbs it.

\step{3}{(3.4a) Receivables} Receivables are part of $TA$ in \eqref{r2ce2-efn}. Holding $L$ and $E$
fixed, $\partial \EFN/\partial AR=1$: in the illustration, raising receivables by $50$ moves EFN from
$-40$ to $-40+50=10$.
\answer{(3.4a)} True: every extra dollar tied up in receivables must be financed.

\step{4}{(3.4b) Cash flow forecasts} A cash flow forecast and a pro forma balance sheet rearrange the
same accounting identities; Higgins' chapter summary states that cash flow forecasts ``yield the same
need for external funding as a pro forma projection, given the same assumptions.''
\answer{(3.4b)} False: with the same assumptions and no errors, the two estimates are identical.

\step{5}{(3.4c) Intra-year needs} \eqref{r2ce2-efn} is evaluated on one date (December 31). A seasonal
firm can build inventory and receivables in mid-year and need financing then, even if by year-end the
cash has come back.
\answer{(3.4c)} False: an annual forecast says nothing about needs between forecast dates; monthly
cash budgets or quarterly pro formas are needed for that.

\intuition EFN is a balancing item, not a forecast of performance. Its sign tells you whether the plan
throws off or absorbs cash \emph{at the forecast date}, and anything that adds to assets raises it
one-for-one.
\pitfalls
\begin{itemize}
  \item Reading a negative EFN as ``the forecast is wrong''; it is a cash surplus.
  \item Believing different forecasting methods give different needs; differences signal an error.
  \item Using a year-end forecast to manage a seasonal business.
\end{itemize}
\bookans 3.1 matches the book (p.~405): ``A negative value implies that the company has excess cash above
its desired minimum'', confirmed by adding it to cash. 3.4: no answer in the book.
\refbib{Higgins Ch.~3, ``Pro Forma Statements'' and chapter summary, pp.~80--102.}

% =============================================================================
\exsec{3. Corpora\c{c}\~ao Felipe: the plug variable --- Ross 4.1 and 4.2, p.~117}

\asked (4.1) Pro forma statements when sales and \emph{every} balance sheet item grow 15\%, and which
variable closes the balance sheet. (4.2) Same forecast, but half of net income is paid as dividends
and debt and equity do not grow with sales: find the external financing needed (NAF).

\approach A pro forma model has one degree of freedom too many once we fix sales growth, the
percentages, and the financing; one item must be the \emph{variable de fechamento} (plug). In 4.1 the
balance sheet is fully specified, so dividends become the plug; in 4.2 the payout is fixed, so the
plug moves to external financing.

\notation
\begin{itemize}
  \item $g=15\%$: sales growth; $S$, $C$, $NI$: sales, costs, net income (no taxes).
  \item $A$, $D$, $E$: assets, debt, equity; $Div$: dividends; $\EFN$ (NAF): external financing needed.
\end{itemize}

\step{1}{Pro forma income statement (both parts)} Costs vary with sales:
\begin{align}
  S_1 &= 23{,}000(1.15) = 26{,}450, \qquad C_1 = 16{,}700(1.15) = 19{,}205, \label{r2ce3-S}\\
  NI_1 &= S_1-C_1 = 26{,}450-19{,}205 = 7{,}245. \label{r2ce3-NI}
\end{align}

\step{2}{(4.1) Balance sheet grown 15\%} Every item grows 15\%:
$A_1=15{,}800(1.15)=18{,}170$, $D_1=5{,}200(1.15)=5{,}980$, $E_1=10{,}600(1.15)=12{,}190$, and
$5{,}980+12{,}190=18{,}170$, so it balances.

\step{3}{(4.1) Find the plug} Equity can only change through retained earnings (no share issues), so
$E_1=E_0+NI_1-Div_1$. Solving for dividends,
\begin{align}
  Div_1 &= NI_1-(E_1-E_0) = 7{,}245-(12{,}190-10{,}600) = 7{,}245-1{,}590 = 5{,}655. \label{r2ce3-div}
\end{align}
If no dividend were paid, equity would be $10{,}600+7{,}245=17{,}845$ and the balance sheet would not
balance; the dividend of $5{,}655$ is what makes it do so.
\answer{(4.1)} Sales $26{,}450$, costs $19{,}205$, net income $7{,}245$; assets $18{,}170$ = debt
$5{,}980$ + equity $12{,}190$. The plug (\emph{vari\'avel de fechamento}) is the dividend, $\$5{,}655$.

\step{4}{(4.2) Fixed payout, fixed debt} Dividends are $0.5(7{,}245)=3{,}622.5$, so retained earnings
add $3{,}622.5$ and
\begin{align}
  E_1 &= 10{,}600+3{,}622.5 = 14{,}222.5, \qquad D_1 = 5{,}200, \nonumber\\
  D_1+E_1 &= 5{,}200+14{,}222.5 = 19{,}422.5, \nonumber\\
  \EFN &= A_1-(D_1+E_1) = 18{,}170-19{,}422.5 = -1{,}252.5. \label{r2ce3-efn}
\end{align}
\answer{(4.2)} NAF $=-\$1{,}252.5$: a surplus. The firm can pay down debt, add to cash, or raise the
dividend by that amount.

\intuition Once growth and the asset intensity are fixed, \emph{something} must give: dividends,
borrowing or new equity. Ross' point is that the pro forma does not tell you which; the manager
chooses the plug, and the model reports the consequence.
\pitfalls
\begin{itemize}
  \item Growing equity by 15\% \emph{and} adding net income: equity changes only through retained earnings.
  \item Reporting a negative NAF as zero; the sign carries information (excess funds).
\end{itemize}
\refbib{Ross \S4.3, ``O modelo da porcentagem de vendas'', pp.~99--104; Exemplo pp.~97--98.}

% =============================================================================
\exsec{4. Internal and sustainable growth rates --- Berk 18.16, p.~561}

\asked (a) The fastest growth the firm can finance with retained earnings alone; (b) the fastest growth
with retained earnings plus the debt that keeps $D/E$ constant; (c) (b) again when 40\% of income is
paid out.

\approach Berk's equations (18.4)--(18.5): internal growth $=$ ROA $\times$ retention, sustainable
growth $=$ ROE $\times$ retention, both on \emph{beginning-of-period} assets and equity. With retention
$b$, equity grows by $b\cdot NI$; assets can grow at that rate relative to assets (no external money)
or relative to equity (debt grows in step).

\notation
\begin{itemize}
  \item $NI=50{,}000$; $A_0=400{,}000$, $E_0=250{,}000$: beginning assets and equity.
  \item $b=1-\text{payout}$: retention (plowback) ratio; $\text{ROA}=NI/A_0$, $\text{ROE}=NI/E_0$.
\end{itemize}

\step{1}{(a) Internal growth} With payout 0\%, $b=1$:
\begin{align}
  g_{int} &= \text{ROA}\times b = \frac{50{,}000}{400{,}000}\times 1 = 0.125 = 12.5\%. \label{r2ce4-int}
\end{align}
\answer{(a)} 12.5\%.

\step{2}{(b) Sustainable growth}
\begin{align}
  g_{sus} &= \text{ROE}\times b = \frac{50{,}000}{250{,}000}\times 1 = 0.20 = 20\%. \label{r2ce4-sus}
\end{align}
\answer{(b)} 20\%.

\step{3}{(c) Payout 40\%} Now $b=1-0.40=0.60$:
\begin{align}
  g_{sus} &= (0.20)(0.60) = 0.12 = 12\%. \label{r2ce4-sus40}
\end{align}
\answer{(c)} 12\%.

\intuition Retained earnings are the only internal source of equity. Measured against assets they bound
growth without any outside money; measured against equity they bound growth when debt is allowed to
grow in proportion. Leverage ($A_0/E_0=1.6$) is why (b) exceeds (a): $20\%=12.5\%\times1.6$.
\pitfalls
\begin{itemize}
  \item Using end-of-year equity or assets with the simple formulas (then the formulas are
        $x b/(1-xb)$; see Exercise 13).
  \item Forgetting to multiply by the retention ratio when there is a payout.
\end{itemize}
\refbib{Berk \S18.4, Eqs.~(18.4)--(18.5), p.~548; Example 18.3, p.~549.}

% =============================================================================
\exsec{5. Universal Products: finishing a pro forma --- Higgins 3.5, p.~104}

\asked (a) Net fixed assets at the end of 2022 given \$500 thousand of purchases. (b) Owners' equity
with a 50\% payout and no equity issues. (c) The resulting external funding required.

\approach Two roll-forward identities fill the blanks: fixed assets evolve by capex less depreciation;
equity by retained earnings. EFN is then the plug of the balance sheet (Exercise 2).

\notation
\begin{itemize}
  \item $NFA$: net fixed assets; $CapEx=500$; $Dep_{2022}=260$ (\$ thousands).
  \item $NI_{2022}=156$; payout $=50\%$; $E$: owners' equity; $CA$, $CL$, $LTD$ as in the table.
\end{itemize}

\step{1}{(a) Fixed assets} Net fixed assets grow by purchases and fall by depreciation:
\begin{align}
  NFA_{2022} &= NFA_{2021}+CapEx-Dep_{2022} = 2{,}000+500-260 = 2{,}240. \label{r2ce5-nfa}
\end{align}
\answer{(a)} \$2,240 thousand.

\step{2}{(b) Equity} Equity grows by retained income only:
\begin{align}
  E_{2022} &= E_{2021}+NI_{2022}(1-0.5) = 2{,}950+156(0.5) = 2{,}950+78 = 3{,}028. \label{r2ce5-eq}
\end{align}
\answer{(b)} \$3,028 thousand.

\step{3}{(c) The plug} Totals and the difference:
\begin{align}
  TA_{2022} &= CA+NFA = 2{,}750+2{,}240 = 4{,}990, \nonumber\\
  (L+E)_{2022} &= CL+LTD+E = 1{,}155+500+3{,}028 = 4{,}683, \nonumber\\
  \EFN_{2022} &= 4{,}990-4{,}683 = 307. \label{r2ce5-efn}
\end{align}
\answer{(c)} External funding required $=\$307$ thousand.

\intuition Assets grow by $4{,}990-4{,}500=490$, but spontaneous liabilities supply only
$1{,}155-1{,}050=105$ and retained earnings $78$; the remaining $490-105-78=307$ must come from outside.
\pitfalls
\begin{itemize}
  \item Adding all of net income to equity (ignoring the 50\% dividend).
  \item Forgetting depreciation in the fixed-asset roll-forward, or subtracting 2021's (250) instead of 2022's (260).
\end{itemize}
\bookans Matches the book (p.~405): $NFA=2{,}240$, equity $3{,}028$, EFN $307$.
\refbib{Higgins Ch.~3, ``Pro Forma Statements'', pp.~80--88.}

% =============================================================================
\exsec{6. Net new financing and how to split it --- Berk 18.2 and 18.3, p.~558}

\asked (18.2) The external financing needed next year given ending assets, net income, payout and the
growth of non-debt liabilities. (18.3) With beginning debt of \$100,000, how much of that financing must
be debt and how much equity to keep the debt-equity ratio unchanged.

\approach Berk's identity: Net New Financing $=$ Projected Assets $-$ Projected Liabilities and Equity,
where projected liabilities and equity include only the \emph{spontaneous} sources (payables) and
retained earnings. Part 18.3 is a financing-policy choice: the total is fixed by 18.2; the policy splits it.

\notation
\begin{itemize}
  \item $TA_1=500{,}000$; $NI=50{,}000$; payout $10\%$; $E_0=300{,}000$; $L_0=120{,}000$ (of which $D_0=100{,}000$ debt in 18.3).
  \item $NNF$: net new financing; $\Delta D$, $\Delta E^{new}$: new debt and new equity issued.
\end{itemize}

\step{1}{(18.2) Projected sources} Retained earnings are $50{,}000(1-0.10)=45{,}000$, so
\begin{align}
  E_1 &= 300{,}000+45{,}000 = 345{,}000, \qquad L_1 = 120{,}000+10{,}000 = 130{,}000, \nonumber\\
  NNF &= TA_1-(L_1+E_1) = 500{,}000-(130{,}000+345{,}000) = 500{,}000-475{,}000 = 25{,}000. \label{r2ce6-nnf}
\end{align}
\answer{(18.2)} Net new financing $=\$25{,}000$.

\step{2}{(18.3) Target capital structure} Non-debt liabilities are $120{,}000-100{,}000=20{,}000$ at the
start and $20{,}000+10{,}000=30{,}000$ at the end. Debt plus equity must therefore fund
$500{,}000-30{,}000=470{,}000$ of assets, in the initial proportion $D/E=100{,}000/300{,}000=1/3$, i.e.\
$D=\tfrac14$ and $E=\tfrac34$ of the total:
\begin{align}
  E_1 &= \tfrac34(470{,}000) = 352{,}500, \qquad D_1 = \tfrac14(470{,}000) = 117{,}500. \label{r2ce6-split}
\end{align}

\step{3}{(18.3) Issues} Equity already reaches $345{,}000$ through retention; debt starts at $100{,}000$:
\begin{align}
  \Delta E^{new} &= 352{,}500-345{,}000 = 7{,}500, \qquad \Delta D = 117{,}500-100{,}000 = 17{,}500, \nonumber
\end{align}
and $7{,}500+17{,}500=25{,}000=NNF$.
\answer{(18.3)} Issue \$7,500 of equity and \$17,500 of debt.

\intuition Net new financing is fixed by the operating plan; how it is raised is a separate capital
structure decision. Keeping $D/E$ constant forces some equity issuance here because retained earnings
($45{,}000$) are too small relative to the asset growth.
\pitfalls
\begin{itemize}
  \item Treating payables as debt when computing $D/E$; they are spontaneous operating liabilities.
  \item Applying $D/E=1/3$ to the \$25,000 alone; the ratio must hold for the \emph{stocks}, and retained
        earnings already added equity.
\end{itemize}
\refbib{Berk \S18.2, ``Pro Forma Balance Sheet'' and ``Making the Balance Sheet Balance'', pp.~539--541.}

% =============================================================================
\exsec{7. Companhia Jo\~ao Herdeiro: percent-of-sales balance sheet --- Ross 4.9--4.11, pp.~118--119}

\asked (4.9) The pro forma income statement at 20\% sales growth with constant payout, and the addition
to retained earnings. (4.10) Each balance sheet item as a percent of sales, ``n/a'' where it does not
vary with sales. (4.11) The pro forma balance sheet at 15\% growth with no new debt or equity, and the NAF.

\approach Percent of sales: costs, assets and payables scale with sales; notes payable, long-term debt
and paid-in capital stay; retained earnings grow by the retained part of the new net income. NAF is the plug.

\notation
\begin{itemize}
  \item $S_0=38{,}000$, $C_0=18{,}400$; tax rate $34\%$; $Div_0=5{,}200$, $NI_0=12{,}936$.
  \item $b=7{,}736/12{,}936=0.598021$: retention ratio; $1-b=0.401979$: payout.
\end{itemize}

\step{1}{(4.9) Income statement at $g=20\%$}
\begin{align}
  S_1 &= 38{,}000(1.2) = 45{,}600, \quad C_1 = 18{,}400(1.2) = 22{,}080, \quad
  EBT_1 = 45{,}600-22{,}080 = 23{,}520, \nonumber\\
  T_1 &= 0.34(23{,}520) = 7{,}996.8, \qquad NI_1 = 23{,}520-7{,}996.8 = 15{,}523.2. \label{r2ce7-ni}
\end{align}
Since all lines scale with sales, $NI_1=1.2(12{,}936)=15{,}523.2$ too. With a constant payout,
$Div_1=15{,}523.2(0.401979)=6{,}240$ and the addition to retained earnings is
$15{,}523.2(0.598021)=9{,}283.2$ $(=1.2\times7{,}736)$.
\answer{(4.9)} Addition to retained earnings $=\$9{,}283.20$ (dividends $\$6{,}240$).

\step{2}{(4.10) Percent of sales} Divide each item by $S_0=38{,}000$:
\begin{center}\small
\begin{tabular}{lrr@{\qquad}lrr}
\toprule
Ativo & \$ & \% sales & Passivo e PL & \$ & \% sales\\
\midrule
Caixa & 3,050 & 8.03 & Contas a pagar & 1,300 & 3.42\\
Contas a receber & 6,900 & 18.16 & Empr\'estimos & 6,800 & n/a\\
Estoque & 7,600 & 20.00 & Total circulante & 8,100 & n/a\\
Total circulante & 17,550 & 46.18 & Passivo n\~ao circulante & 25,000 & n/a\\
Instala\c{c}\~oes e equip. & 34,500 & 90.79 & Capital social & 15,000 & n/a\\
 & & & Lucros retidos & 3,950 & n/a\\
 & & & Total PL & 18,950 & n/a\\
Total do ativo & 52,050 & 136.97 & Total & 52,050 & n/a\\
\bottomrule
\end{tabular}
\end{center}
The current-liabilities total is n/a because only part of it (payables) varies with sales.
\answer{(4.10)} See the table: all assets and payables scale (total assets $136.97\%$ of sales); loans,
long-term debt and equity are n/a.

\step{3}{(4.11) Income at $g=15\%$} $S_1=43{,}700$, $C_1=21{,}160$, $EBT_1=22{,}540$,
$T_1=0.34(22{,}540)=7{,}663.6$, $NI_1=22{,}540-7{,}663.6=14{,}876.4$; dividends
$14{,}876.4(0.401979)=5{,}980$; addition to retained earnings $14{,}876.4(0.598021)=8{,}896.4$.

\step{4}{(4.11) Pro forma balance sheet} Scaling assets and payables by 1.15 and rolling retained
earnings forward ($3{,}950+8{,}896.4=12{,}846.4$):
\begin{center}\small
\begin{tabular}{lr@{\qquad}lr}
\toprule
Ativo & \$ & Passivo e PL & \$\\
\midrule
Caixa & 3,507.5 & Contas a pagar & 1,495.0\\
Contas a receber & 7,935.0 & Empr\'estimos & 6,800.0\\
Estoque & 8,740.0 & Total circulante & 8,295.0\\
Total circulante & 20,182.5 & Passivo n\~ao circulante & 25,000.0\\
Instala\c{c}\~oes e equip. & 39,675.0 & Capital social & 15,000.0\\
 & & Lucros retidos & 12,846.4\\
 & & Total PL & 27,846.4\\
Total do ativo & 59,857.5 & Total & 61,141.4\\
\bottomrule
\end{tabular}
\end{center}
\begin{align}
  \text{NAF} &= 59{,}857.5-(8{,}295+25{,}000+27{,}846.4) = 59{,}857.5-61{,}141.4 = -1{,}283.9. \label{r2ce7-naf}
\end{align}
\answer{(4.11)} NAF $=-\$1{,}283.90$: at 15\% growth the firm generates a surplus.

\intuition Asset needs rise by $0.15\times52{,}050=7{,}807.5$; payables supply $195$ and retained earnings
$8{,}896.4$, more than enough. A very profitable firm (net margin 34\%) can finance moderate growth
internally; the surplus would turn into a deficit at a higher growth rate (Exercise 14 shows how to find
the switch point).
\pitfalls
\begin{itemize}
  \item Using the 20\% income statement of 4.9 in 4.11; the growth rate changes to 15\%.
  \item Scaling loans (\emph{empr\'estimos}) with sales: they are a financing decision.
\end{itemize}
\refbib{Ross \S4.3, pp.~99--104.}

% =============================================================================
\exsec{8. Global Corp.: two-pass pro forma --- Berk 18.8--18.11, p.~560}

\asked (18.8) Next year's costs, depreciation, net income and the operating balance sheet items at 8\%
sales growth. (18.9) Equity with a 50\% payout. (18.10) Net new financing. (18.11) The payout that keeps
net new financing at or below \$9 million.

\approach Berk's KMS template (Tables 18.2--18.3): costs and depreciation as a percent of sales,
interest expense unchanged, tax at the current rate; assets and payables as a percent of sales; debt
unchanged in the first pass; equity rolls forward by retained earnings; net new financing is the gap.
Part 18.11 inverts the plug: fix NNF, solve for retention.

\notation
\begin{itemize}
  \item $g=8\%$; amounts in \$ millions; $\tau=0.7/2.7=0.259259$: Global's tax rate (taxes/pretax income).
  \item $p$: payout; $NNF$: net new financing.
\end{itemize}

\step{1}{(18.8) Income statement} Every operating line scales by 1.08; interest stays at 7.7:
\begin{align}
  S_1 &= 186.7(1.08) = 201.636, \quad C_1 = 175.1(1.08) = 189.108, \quad Dep_1 = 1.2(1.08) = 1.296, \nonumber\\
  EBIT_1 &= 201.636-189.108-1.296 = 11.232, \qquad EBT_1 = 11.232-7.7 = 3.532, \nonumber\\
  T_1 &= 0.259259(3.532) = 0.915704, \qquad NI_1 = 3.532-0.915704 = 2.616296. \label{r2ce8-ni}
\end{align}
(Depreciation is $1.2/186.7=0.6427\%$ of sales; scaling by 1.08 is the same thing.)

\step{2}{(18.8) Balance sheet items} Each scales by 1.08: cash $23.2(1.08)=25.056$; receivables
$18.5(1.08)=19.98$; inventory $15.3(1.08)=16.524$ (current assets $61.56$); PP\&E $113.1(1.08)=122.148$;
total assets $183.708$; payables $34.7(1.08)=37.476$.
\answer{(18.8)} (a) 189.108; (b) 1.296; (c) 2.616; (d) 25.056; (e) 19.98; (f) 16.524; (g) 122.148; (h) 37.476 (\$ millions).

\step{3}{(18.9) Equity} Retained earnings are $0.5(2.616296)=1.308148$:
\begin{align}
  E_1 &= 22.2+1.308148 = 23.508148. \label{r2ce8-eq}
\end{align}
\answer{(18.9)} Stockholders' equity $\approx\$23.51$ million.

\step{4}{(18.10) Net new financing} Debt stays at $113.2$ in the first pass:
\begin{align}
  (L+E)_1 &= 37.476+113.2+23.508148 = 174.184148, \nonumber\\
  NNF &= 183.708-174.184148 = 9.523852. \label{r2ce8-nnf}
\end{align}
\answer{(18.10)} Net new financing $\approx\$9.52$ million.

\step{5}{(18.11) Payout cap} Write NNF as a function of the payout:
$NNF(p)=183.708-37.476-113.2-22.2-(1-p)NI_1=10.832-(1-p)(2.616296)$. Requiring $NNF\le 9$:
\begin{align}
  (1-p)(2.616296) &\ge 10.832-9 = 1.832, \nonumber\\
  1-p &\ge \frac{1.832}{2.616296} = 0.700227, \nonumber\\
  p &\le 0.299773. \label{r2ce8-p}
\end{align}
Dividends fall from $1.308$ to at most $2.616296-1.832=0.784$ million, i.e.\ by the $0.524$ excess.
\answer{(18.11)} Global must cut its payout from 50\% to at most about 30\% (dividends $\le\$0.78$ million).

\intuition Global is highly levered with a thin margin (net income \$2.0 million on \$186.7 million of
sales), so 8\% growth in a \$170 million asset base creates a financing need ($13.6$ of new assets) that
retained earnings barely dent. The payout is the only internal lever left; cutting it from 50\% to 30\%
saves just \$0.52 million.
\pitfalls
\begin{itemize}
  \item Growing interest expense with sales; the method holds it fixed until the financing is decided.
  \item Using a 35\% tax rate (the KMS rate) instead of Global's own $0.7/2.7$.
  \item Solving 18.11 by scaling dividends instead of solving the inequality for the payout.
\end{itemize}
\refbib{Berk \S18.2, Tables 18.1--18.4 and Example 18.2, pp.~537--541.}

% =============================================================================
\exsec{9. KMS growth vs.\ its SGR; new debt; D/E at the internal rate --- Berk 18.17--18.19, p.~561}

\asked (18.17) Whether the KMS plan grows sales faster or slower than its sustainable growth rate (SGR).
(18.18) How much debt a firm must issue if it grows exactly at its SGR. (18.19) IZAX's new $D/E$ if its
assets grow at its internal growth rate.

\approach SGR $=$ ROE $\times b$ (beginning equity), Eq.~(18.5): the growth that leaves $D/E$ constant
with no new equity. Comparing planned growth with SGR tells us whether leverage, payout or equity issues
must change. At the internal rate no external money is raised, so debt is frozen and $D/E$ falls.

\notation
\begin{itemize}
  \item $g_t=S_t/S_{t-1}-1$: KMS sales growth; $\text{SGR}_t=(NI_t-Div_t)/E_{t-1}$; $\text{ROE}_t=NI_t/E_{t-1}$.
  \item 18.18: $\text{ROE}=12\%$, payout $25\%$, $E_0=600{,}000$, $D_0=400{,}000$.
  \item 18.19: $A_0=400{,}000$, $D_0=100{,}000$, $E_0=300{,}000$, $NI=20{,}000$, $Div=5{,}000$.
\end{itemize}

\step{1}{(18.17) Growth and SGR year by year} Using Tables 18.8 and 18.17 (\$000s); e.g.\ 2011:
$g=88{,}369/74{,}889-1=18.00\%$ and $\text{SGR}=(7{,}600-4{,}786)/74{,}134=2{,}814/74{,}134=3.80\%$.
\begin{center}\small
\begin{tabular}{lrrrrr}
\toprule
 & 2011 & 2012 & 2013 & 2014 & 2015\\
\midrule
Sales growth $g_t$ (\%) & 18.00 & 16.84 & 16.03 & 15.34 & 14.75\\
SGR with the plan's dividends (\%) & 3.80 & 8.19 & 8.29 & 8.40 & 8.51\\
ROE $=$ SGR at full retention (\%) & 10.25 & 11.45 & 13.38 & 15.24 & 17.01\\
\bottomrule
\end{tabular}
\end{center}
Even under the 2010 policy (30\% payout, Section 18.2) the SGR was $6{,}941/69{,}275\times0.7=7.01\%$.
\answer{(18.17)} Faster: sales grow 15--18\% a year against an SGR of 4--9\%. KMS covers the gap by
raising leverage --- book $D/E$ jumps from $4{,}500/74{,}134=6.07\%$ to $24{,}500/76{,}948=31.84\%$ in 2011 ---
and, through 2014, growth exceeds even the full-retention ROE.

\step{2}{(18.18) Debt at the SGR}
\begin{align}
  \text{SGR} &= 0.12(1-0.25) = 0.09, \qquad \Delta E = 600{,}000(0.12)(0.75) = 54{,}000 = 9\%\ \text{of } E_0, \nonumber\\
  \Delta D &= 0.09(400{,}000) = 36{,}000. \label{r2ce9-dd}
\end{align}
Check: $(400{,}000+36{,}000)/(600{,}000+54{,}000)=436{,}000/654{,}000=2/3=400{,}000/600{,}000$.
\answer{(18.18)} Issue \$36,000 of new debt.

\step{3}{(18.19) IZAX at the internal rate}
\begin{align}
  \text{ROA} &= 20{,}000/400{,}000 = 0.05, \quad b = 15{,}000/20{,}000 = 0.75, \quad g_{int} = (0.05)(0.75) = 0.0375, \nonumber\\
  A_1 &= 400{,}000(1.0375) = 415{,}000, \quad E_1 = 300{,}000+15{,}000 = 315{,}000, \quad D_1 = 100{,}000, \nonumber\\
  D_1/E_1 &= 100{,}000/315{,}000 = 0.3175. \label{r2ce9-de}
\end{align}
\answer{(18.19)} $D/E=0.3175$ (down from $0.3333$).

\intuition The SGR is a benchmark, not a target. KMS's planned growth is far above it, which only says
that the plan needs outside money (here a \$20 million bond); whether that growth is \emph{worth} it
is an NPV question (Exercise 15). At the internal rate equity grows and debt does not, so leverage
drifts down; at the sustainable rate debt grows with equity and leverage stays put.
\pitfalls
\begin{itemize}
  \item Calling growth above the SGR ``bad'': Berk stresses that only NPV decides value.
  \item Computing ROE on end-of-year equity with the simple formula.
\end{itemize}
\refbib{Berk \S18.4, pp.~547--551; Tables 18.8, 18.17, pp.~545, 562.}

% =============================================================================
\exsec{10. Under Armour: sustainable growth and its sensitivity --- Higgins 4.8, pp.~135--136}

\asked (a) The sustainable growth rate each year 2009--2013. (b) Whether actual growth outran it.
(c) How the company coped. (d) The 2013 SGR with an 8\% margin, with leverage of 2.40, and with both.

\approach Higgins' $g^*=PRA\hat T$: profit margin $\times$ retention $\times$ asset turnover $\times$
assets/beginning equity. A firm growing faster than $g^*$ must improve operations ($P$, $A$) or change
financial policy ($R$, $\hat T$) or issue equity. Part (d) is a one-at-a-time sensitivity analysis.

\notation
\begin{itemize}
  \item $P$ (\%), $R=1$, $A$ (times), $\hat T$ (times): the four ratios in the table; $g$: actual sales growth.
\end{itemize}

\step{1}{(a) The product of four ratios} With $R=1$ in every year, $g^*=P\cdot A\cdot\hat T$:
\begin{align}
  g^*_{2009} &= 5.5(1.6)(1.7) = 14.96\%, & g^*_{2010} &= 6.4(1.6)(1.7) = 17.41\%, \nonumber\\
  g^*_{2011} &= 6.6(1.6)(1.9) = 20.06\%, & g^*_{2012} &= 7.0(1.6)(1.8) = 20.16\%, \nonumber\\
  g^*_{2013} &= 7.0(1.5)(1.9) = 19.95\%. & & \label{r2ce10-g}
\end{align}
\answer{(a)} 14.96\%, 17.41\%, 20.06\%, 20.16\%, 19.95\%.

\step{2}{(b) Compare with actual growth} Actual growth was 18.1, 24.2, 38.4, 24.6 and 27.1\%: above
$g^*$ every year, by up to $38.4-20.06=18.34$ points in 2011.
\answer{(b)} Yes: a too-much-growth challenge --- growing faster than equity could be built internally.

\step{3}{(c) Which ratios moved} Retention was already 100\%. The margin rose (5.5\% $\to$ 7.0\%) and
leverage rose ($\hat T$: 1.7 $\to$ 1.9); turnover slipped (1.6 $\to$ 1.5), which worked the other way.
Because each year's $g^*$ uses that year's ratios, growth above $g^*$ in a year must have been met by
ratios rising further (or new equity) during that year.
\answer{(c)} By raising profitability and, above all, financial leverage (more debt per dollar of
beginning equity), with no dividends; operating efficiency (turnover) did not help.

\step{4}{(d) Sensitivity of 2013 $g^*$}
\begin{align}
  P=8.0\%:&\quad g^* = 8.0(1.5)(1.9) = 22.8\%, \nonumber\\
  \hat T=2.40:&\quad g^* = 7.0(1.5)(2.4) = 25.2\%, \nonumber\\
  \text{both}:&\quad g^* = 8.0(1.5)(2.4) = 28.8\%. \label{r2ce10-d}
\end{align}
\answer{(d)} 22.8\%, 25.2\% and 28.8\%; only the combination exceeds 2013's actual 27.1\%.

\intuition $g^*$ is linear in each lever, so the levers multiply: a 14\% rise in the margin and a 26\%
rise in leverage together lift $g^*$ by 44\%. Leverage is the easiest lever to pull and the riskiest
to rely on.
\pitfalls
\begin{itemize}
  \item Adding the percentage effects in (d) instead of multiplying the ratios.
  \item Using end-of-period assets/equity; Higgins' $\hat T$ uses beginning equity.
\end{itemize}
\refbib{Higgins Ch.~4, ``Sustainable Growth'' and ``What If Actual Growth Exceeds Sustainable Growth?'', pp.~112--123.}

% =============================================================================
\exsec{11. Turismo da On\c{c}a: NAF, capacity, constant D/E --- Ross 4.25--4.27, p.~121}

\asked (4.25) The NAF for 20\% growth at full capacity with no new debt or equity, interest, tax rate
and payout constant. (4.26) The NAF if the firm ran at 80\% of capacity. (4.27) The NAF if the firm
keeps its debt-equity ratio constant.

\approach Percent of sales for costs, other expenses, current assets and payables; fixed assets scale
only if there is no idle capacity; interest is \emph{fixed}. In 4.27 the financing policy allows debt
to grow with equity, so part of the gap is filled by planned borrowing.

\notation
\begin{itemize}
  \item $S_0=929{,}000$; $\tau=60{,}550/173{,}000=0.35$; payout $=33{,}735/112{,}450=0.30$; interest $14{,}000$.
  \item ``Total debt'' (Ross) = all liabilities: $D_0=85{,}000+158{,}000=243{,}000$; $E_0=322{,}900$.
\end{itemize}

\step{1}{Income statement at $g=20\%$}
\begin{align}
  S_1 &= 1{,}114{,}800, \quad \text{Custos} = 867{,}600, \quad \text{Outras} = 22{,}800, \quad EBIT = 224{,}400, \nonumber\\
  EBT &= 224{,}400-14{,}000 = 210{,}400, \quad T = 0.35(210{,}400) = 73{,}640, \quad NI = 136{,}760, \nonumber\\
  Div &= 0.30(136{,}760) = 41{,}028, \qquad \text{Acr\'escimo aos lucros retidos} = 95{,}732. \label{r2ce11-is}
\end{align}

\step{2}{(4.25) Balance sheet at full capacity} Current assets $152{,}900(1.2)=183{,}480$
(cash $30{,}360$, receivables $48{,}840$, inventory $104{,}280$); fixed assets $413{,}000(1.2)=495{,}600$;
total $679{,}080$. Payables $81{,}600$; notes $17{,}000$; long-term $158{,}000$; equity
$140{,}000+182{,}900+95{,}732=418{,}632$; total $81{,}600+17{,}000+158{,}000+418{,}632=675{,}232$.
\begin{align}
  \text{NAF} &= 679{,}080-675{,}232 = 3{,}848. \label{r2ce11-25}
\end{align}
\answer{(4.25)} NAF $=\$3{,}848$.

\step{3}{(4.26) Idle capacity} Full-capacity sales are $929{,}000/0.80=1{,}161{,}250>1{,}114{,}800$, so no new
fixed assets are needed (utilization rises to 96\%). Total assets $=183{,}480+413{,}000=596{,}480$:
\begin{align}
  \text{NAF} &= 596{,}480-675{,}232 = -78{,}752. \label{r2ce11-26}
\end{align}
\answer{(4.26)} NAF $=-\$78{,}752$ (a surplus).

\step{4}{(4.27) Constant total-debt/equity} $D/E=243{,}000/322{,}900=0.752555$. With equity
$418{,}632$, permitted total debt is $0.752555(418{,}632)=315{,}043.59$; total sources
$418{,}632+315{,}043.59=733{,}675.59$:
\begin{align}
  \text{NAF} &= 679{,}080-733{,}675.59 = -54{,}595.59. \label{r2ce11-27}
\end{align}
Existing and spontaneous liabilities total $81{,}600+17{,}000+158{,}000=256{,}600$, so the policy allows
$315{,}043.59-256{,}600=58{,}443.59$ of new borrowing, which more than covers the $3{,}848$ gap.
\answer{(4.27)} NAF $\approx-\$54{,}596$: the debt the policy allows exceeds the need by that amount.

\intuition Three different answers from one forecast: the NAF depends as much on assumptions about
capacity and financing policy as on the sales forecast. Idle capacity is worth \$82,600 of financing here
($495{,}600-413{,}000$).
\pitfalls
\begin{itemize}
  \item Growing interest expense with sales (the statement holds it constant).
  \item In 4.26, cutting fixed assets below the current level; they stay at 413,000.
  \item In 4.27, applying $D/E$ only to interest-bearing debt; Ross' ratio uses total liabilities (the
        alternative definition changes the number, so state it).
\end{itemize}
\refbib{Ross \S4.3, ``Uma situa\c{c}\~ao alternativa'', pp.~103--104; \S4.4, pp.~106--111.}

% =============================================================================
\exsec{12. R\&E Supplies: extending the spreadsheet, sensitivity and scenario --- Higgins 3.9 (3.8), p.~105}

\asked (3.8, Part 0) The 2023 spreadsheet formulas for net sales and equity. (a) The 2023 external
financing required (EFR) and its comparison with 2022. (b) EFR if COGS falls to 84\% of sales.
(c) EFR in a recession: sales $-5\%$, COGS 88\%, current assets 35\% of sales.

\approach Extend Table 3.5 one column: 2023 sales grow from 2022 forecast sales; 2023 equity starts
from 2022 forecast equity. Interest is charged on long-term debt, its current portion \emph{and} the
EFR itself, which makes EFR circular; we solve the circularity in closed form instead of iterating.
(b) changes one input (sensitivity); (c) changes several together (scenario).

\notation
\begin{itemize}
  \item $X$: 2023 EFR (\$ thousands); $i=10\%$; $\tau=25\%$; payout $50\%$; $LTD=560$, $CPLTD=100$, $NFA=270$.
  \item $S$, $GP$, $SGA$, $CA$, $CL$: 2023 sales, gross profit, SG\&A (11\%), current assets, current liabilities (14.4\%).
  \item $S_{22}=20{,}613(1.25)=25{,}766.25$ and $E_{22}=1{,}837.81$: Table 3.5's unrounded 2022 values (the table prints 25,766 and 1,838).
\end{itemize}

\step{0}{(3.8) Formulas} \texttt{D19 = C19 + C19*D4} (2023 sales from 2022 forecast sales) and
\texttt{D37 = C37 + D28} (2023 equity = 2022 forecast equity + 2023 additions to retained earnings).
\answer{(3.8)} \texttt{D19 = C19 + C19*D4}; \texttt{D37 = C37 + D28}.

\step{1}{Closed form for the circular interest} EFR $=TA-CL-LTD-E$, with
$E=E_{22}+(1-0.5)(1-\tau)\big[GP-SGA-i(LTD+CPLTD+X)\big]$. Let $k=(0.5)(0.75)=0.375$. Then
\begin{align}
  X &= TA-CL-LTD-E_{22}-k\big[GP-SGA-i(660)\big]+k\,i\,X, \nonumber\\
  X(1-k\,i) &= \underbrace{TA-CL-LTD-E_{22}-k\big[GP-SGA-66\big]}_{K}, \qquad X=\frac{K}{1-0.0375}=\frac{K}{0.9625}. \label{r2ce12-X}
\end{align}

\step{2}{(a) Base case} $S=25{,}766.25(1.30)=33{,}496.13$; $COGS=0.86S=28{,}806.67$; $GP=4{,}689.46$;
$SGA=0.11S=3{,}684.57$; $CA=0.29S=9{,}713.88$; $TA=9{,}713.88+270=9{,}983.88$; $CL=0.144S=4{,}823.44$.
\begin{align}
  K &= 9{,}983.88-4{,}823.44-560-1{,}837.81-0.375(4{,}689.46-3{,}684.57-66) = 2{,}410.54, \nonumber\\
  X &= 2{,}410.54/0.9625 = 2{,}504.46. \label{r2ce12-a}
\end{align}
Then interest $=0.1(660+2{,}504.46)=316.45$; EBT $=4{,}689.46-3{,}684.57-316.45=688.44$; tax $172.11$;
EAT $516.33$; dividends $258.16$; equity $1{,}837.81+258.16=2{,}095.97$; liabilities and equity
$4{,}823.44+560+2{,}095.97=7{,}479.42$; and $9{,}983.88-7{,}479.42=2{,}504.46=X$, as it must.
\answer{(a)} EFR $\approx\$2{,}504$ thousand in 2023, about \$986 thousand more than 2022's \$1,518 thousand.

\step{3}{(b) Sensitivity: COGS 84\%} Only $GP$ changes: $GP=0.16S=5{,}359.38$, so
$K=2{,}410.54-0.375(5{,}359.38-4{,}689.46)=2{,}159.32$ and $X=2{,}159.32/0.9625=2{,}243.45$.
\answer{(b)} EFR $\approx\$2{,}243$ thousand: $-\$261$ thousand, or $-10.4\%$.

\step{4}{(c) Recession scenario} $S=25{,}766.25(0.95)=24{,}477.94$; $COGS=0.88S=21{,}540.59$; $GP=2{,}937.35$;
$SGA=2{,}692.57$; $CA=0.35S=8{,}567.28$; $TA=8{,}837.28$; $CL=3{,}524.82$.
\begin{align}
  K &= 8{,}837.28-3{,}524.82-560-1{,}837.81-0.375(2{,}937.35-2{,}692.57-66) = 2{,}847.60, \nonumber\\
  X &= 2{,}847.60/0.9625 = 2{,}958.55. \label{r2ce12-c}
\end{align}
Interest $=361.85$, EBT $=-117.08$, EAT $=-87.81$, and the spreadsheet's dividend formula
\texttt{=C11*C26} pays $-43.90$ (a negative dividend).
\answer{(c)} EFR $\approx\$2{,}959$ thousand, $+18.1\%$ over the base case. If dividends are set to zero
when earnings are negative (the realistic policy), EFR is $\approx\$3{,}006$ thousand.

\intuition The financing need is far more sensitive to working capital than to margins: the recession
\emph{lowers} sales but raises EFR because current assets jump to 35\% of a smaller sales figure and
profits vanish. Sensitivity analysis shows which inputs drive the answer; scenario analysis shows how
they move together.
\pitfalls
\begin{itemize}
  \item Starting 2023 from 2021 data (sales \$20,613, equity \$1,730) instead of the 2022 forecast.
  \item Ignoring the interest on the EFR itself (understates EFR by about $0.0375X\approx\$94$ thousand).
  \item Copying a formula that pays negative dividends without noticing.
\end{itemize}
\bookans Matches the book (p.~406): \$2.504 million; \$2.243 million ($-10.4\%$); \$2.959 million
($+18.2\%$ in the book; $+18.1\%$ here from unrounded values). The book's recession figure uses the
spreadsheet's mechanical (negative) dividend.
\refbib{Higgins Ch.~3, ``Forecasting with Spreadsheets'' and ``Coping with Uncertainty'', pp.~89--98.}

% =============================================================================
\exsec{13. Deriving NAF and the growth formulas --- Ross 4.31--4.33, p.~122}

\asked (4.31) Show $\text{NAF}=-ML(V)b+(A-ML(V)b)g$. (4.32) Use it to recover the internal and
sustainable growth formulas [4.2]--[4.3]. (4.33) Show that the end-of-period formula
$\text{ROE}\,b/(1-\text{ROE}\,b)$ equals $\text{ROE}_{bop}\times b$, and similarly for ROA.

\approach The NAF is uses minus internal sources. All assets scale with sales (asset need $Ag$), the
margin is constant, and no liability is spontaneous; retained earnings are next year's income times $b$.
The sustainable case adds the debt that keeps $D/P$ constant.

\notation
\begin{itemize}
  \item $V$: last year's sales; $A$, $D$, $P$: last year's total assets, debt, equity ($A=D+P$).
  \item $g$: sales growth; $ML$: profit margin; $b$: retention; $\text{ROA}=ML\cdot V/A$, $\text{ROE}=ML\cdot V/P$
        (Ross' definitions, end of the base year).
\end{itemize}

\step{1}{(4.31) Uses and sources} New assets are $A g$; next year's net income is $ML\cdot V(1+g)$, of
which $b$ is retained. So
\begin{align}
  \text{NAF} &= A g - ML(V)b(1+g) \nonumber\\
  &= A g - ML(V)b - ML(V)b\,g \nonumber\\
  &= -ML(V)b + \big(A-ML(V)b\big)g. \label{r2ce13-naf}
\end{align}
\answer{(4.31)} Shown in \eqref{r2ce13-naf}.

\step{2}{(4.32) Internal growth} Set $\text{NAF}=0$ in \eqref{r2ce13-naf}:
\begin{align}
  \big(A-ML(V)b\big)g &= ML(V)b \;\Longrightarrow\; g = \frac{ML(V)b}{A-ML(V)b}. \nonumber
\end{align}
Divide numerator and denominator by $A$ and use $\text{ROA}=ML\cdot V/A$:
\begin{align}
  g_{int} &= \frac{(ML\cdot V/A)\,b}{1-(ML\cdot V/A)\,b} = \frac{\text{ROA}\times b}{1-\text{ROA}\times b}. \label{r2ce13-int}
\end{align}

\step{3}{(4.32) Sustainable growth} Keeping $D/P$ constant lets debt rise by $(D/P)$ per dollar of
retained earnings, so the sources are $(1+D/P)\,ML(V)b(1+g)$ and $1+D/P=A/P$. Setting uses equal to
sources:
\begin{align}
  A g &= \frac{A}{P}ML(V)b(1+g) \;\Longrightarrow\; g\Big(A-\frac{A}{P}ML(V)b\Big) = \frac{A}{P}ML(V)b, \nonumber\\
  g_{sus} &= \frac{(ML\cdot V/P)\,b}{1-(ML\cdot V/P)\,b} = \frac{\text{ROE}\times b}{1-\text{ROE}\times b} \quad(\text{dividing by } A). \label{r2ce13-sus}
\end{align}
Illustration (Ross' Hoffman example: $V=A=500$, $NI=66$, $P=250$, $b=2/3$): $g_{int}=9.65\%$, $g_{sus}=21.36\%$.
\answer{(4.32)} Both chapter formulas [4.2] and [4.3] follow from \eqref{r2ce13-naf}.

\step{4}{(4.33) Beginning- vs.\ end-of-period ROE} With no new shares, end equity is
$P_e=P_b+b\,NI$. Let $\text{ROE}_e=NI/P_e$. Then
\begin{align}
  \frac{\text{ROE}_e\,b}{1-\text{ROE}_e\,b} &= \frac{b\,NI/P_e}{1-b\,NI/P_e} = \frac{b\,NI}{P_e-b\,NI}
  = \frac{b\,NI}{P_b} = \text{ROE}_{b}\times b. \label{r2ce13-roe}
\end{align}
The same steps with assets work when assets grow only by retained earnings (no external financing,
which is exactly the internal-growth case), $A_e=A_b+b\,NI$:
$\text{ROA}_e b/(1-\text{ROA}_e b)=b\,NI/(A_e-b\,NI)=b\,NI/A_b=\text{ROA}_b\times b$. Ross' numbers:
$NI=20$, $b=0.6$, $P_b=80$ give $0.6(20/80)=15\%$ and $0.6(0.2174)/(1-0.6(0.2174))=15\%$; with $A_b=100$, $12\%$.
\answer{(4.33)} \eqref{r2ce13-roe} and its ROA analogue.

\intuition The ``$1-\text{ROE}\,b$'' denominator is just the conversion from end-of-period to
beginning-of-period equity: retained earnings are in the end balance but cannot have financed
themselves. Use one consistent definition and you get the same growth rate.
\pitfalls
\begin{itemize}
  \item Forgetting the $(1+g)$ on next year's earnings, which drops the $-ML(V)b\,g$ term.
  \item Mixing definitions: end-of-period ROE in the simple $\text{ROE}\times b$ understates growth
        ($0.6\times0.2174=13.0\%$ instead of $15\%$ in Ross' example); beginning-of-period ROE in [4.3]
        overstates it ($0.15/(1-0.15)=17.6\%$).
\end{itemize}
\refbib{Ross \S4.4, pp.~106--113, Tabela 4.9.}

% =============================================================================
\exsec{14. Turismo da On\c{c}a: NAF as a function of growth --- Ross 4.28 and 4.29, p.~121}

\asked (4.28) Redo 4.25 at 15\% and 25\% growth, plot NAF against $g$, find the $g$ where NAF $=0$, and
explain why it differs from the internal growth formula. (4.29) The same for 4.27 at 20, 30 and 35\%.

\approach Write NAF as an explicit function of $g$ using the model of Exercise 11: assets and payables
scale with $1+g$; interest is fixed at 14,000; retained earnings
$=(1-0.30)(1-0.35)\big[187{,}000(1+g)-14{,}000\big]=0.455\big[187{,}000(1+g)-14{,}000\big]$.
NAF is then linear in $g$, and its root is the model's own internal (4.28) or sustainable (4.29) rate.

\notation
\begin{itemize}
  \item $RE(g)$: addition to retained earnings at growth $g$; $D/E=0.752555$ (total liabilities/equity).
\end{itemize}

\step{1}{(4.28) NAF$(g)$ with no new debt or equity}
\begin{align}
  \text{NAF}(g) &= 565{,}900(1+g)-\big[68{,}000(1+g)+17{,}000+158{,}000+322{,}900+RE(g)\big] \nonumber\\
  &= (565{,}900-68{,}000-0.455\cdot187{,}000)(1+g)-(497{,}900-0.455\cdot14{,}000) \nonumber\\
  &= 412{,}815(1+g)-491{,}530. \label{r2ce14-28}
\end{align}
At $g=15\%$: $RE=91{,}477.75$, NAF $=-16{,}792.75$; at $20\%$: $RE=95{,}732$, NAF $=3{,}848$; at $25\%$:
$RE=99{,}986.25$, NAF $=24{,}488.75$. Root: $1+g=491{,}530/412{,}815$, so $g=19.07\%$.

\step{2}{(4.28) The textbook formula} $\text{ROA}=112{,}450/565{,}900=19.87\%$, $\text{ROA}\cdot b=0.139097$,
$g_{int}=0.139097/(1-0.139097)=16.16\%$. The formula equals the root of
$565{,}900g-78{,}715(1+g)=0$, i.e.\ it assumes (i) no spontaneous liabilities and (ii) a constant net
margin. Adding (i) back (payables grow with sales) moves the root to $78{,}715/(565{,}900-68{,}000-78{,}715)=18.78\%$;
adding (ii) (fixed interest makes the margin rise with sales) moves it to $19.07\%$.
\answer{(4.28)} NAF $=-16{,}793$; $3{,}848$; $24{,}489$. NAF $=0$ at $g\approx19.07\%$, above the formula's
16.16\% because payables finance part of the growth and fixed interest raises the margin.

\step{3}{(4.29) NAF$(g)$ with constant $D/E$} Sources are equity times $(1+D/E)$:
\begin{align}
  \text{NAF}(g) &= 565{,}900(1+g)-1.752555\big[322{,}900+RE(g)\big] \nonumber\\
  &= \big(565{,}900-1.752555\cdot85{,}085\big)(1+g)-1.752555(322{,}900-6{,}370) \nonumber\\
  &= 416{,}783.86(1+g)-554{,}736.22. \label{r2ce14-29}
\end{align}
At $20\%$: NAF $=-54{,}595.59$; at $30\%$ ($RE=104{,}240.5$): $-12{,}917.21$; at $35\%$
($RE=108{,}494.75$): $7{,}921.99$. Root: $g=554{,}736.22/416{,}783.86-1=33.10\%$.

\step{4}{(4.29) The textbook formula} $\text{ROE}=112{,}450/322{,}900=34.83\%$, $\text{ROE}\cdot b=0.243775$,
$g_{sus}=0.243775/(1-0.243775)=32.24\%$, the root of $565{,}900g=1.752555\times78{,}715(1+g)$
$=137{,}952.36(1+g)$, which assumes a constant margin. Payables do not matter here (they count as debt
in $D/E$); the gap is the fixed interest only.
\answer{(4.29)} NAF $=-54{,}596$; $-12{,}917$; $7{,}922$. NAF $=0$ at $g\approx33.10\%$, above the formula's
32.24\% because interest does not grow with sales, so the margin rises with growth.

\begin{center}
\begin{tikzpicture}[x=0.32cm,y=0.032cm]
  \draw[->] (10,0) -- (41,0) node[right]{\small $g$ (\%)};
  \draw[->] (10,-100) -- (10,95) node[above]{\small NAF (\$000)};
  \foreach \x in {15,20,25,30,35,40} \draw (\x,2)--(\x,-2) node[below]{\scriptsize \x};
  \foreach \y in {-75,-50,-25,25,50,75} \draw (9.7,\y)--(10.3,\y) node[left=2pt]{\scriptsize \y};
  % 4.28: 412.815(1+g) - 491.53 (thousands); 4.29: 416.784(1+g) - 554.736
  \draw[thick,blue!70!black] (10,-37.4335) -- (40,86.411) node[right]{\scriptsize 4.28};
  \draw[thick,red!70!black] (10,-96.2741) -- (40,28.7616) node[right]{\scriptsize 4.29};
  \foreach \p in {(15,-16.79275),(20,3.848),(25,24.48875)} \fill[blue!70!black] \p circle (2pt);
  \foreach \p in {(20,-54.5956),(30,-12.9172),(35,7.9220)} \fill[red!70!black] \p circle (2pt);
  \draw[dashed] (19.07,0) -- (19.07,-12) node[below]{\scriptsize 19.07};
  \draw[dashed] (33.10,0) -- (33.10,-30) node[below]{\scriptsize 33.10};
\end{tikzpicture}
\end{center}

\intuition NAF rises linearly with growth: each extra point of growth needs $\$4{,}128$ more financing
in 4.28 ($412{,}815/100$) and $\$4{,}168$ in 4.29. The closed-form growth rates are exact only in the
stylized world where every cost scales with sales; real pro formas have fixed costs (here interest), so
always verify with the model itself.
\pitfalls
\begin{itemize}
  \item Treating the formula rates as the model's break-even rates without checking its assumptions.
  \item Plotting NAF against $g$ with curvature: in this model it is a straight line.
\end{itemize}
\refbib{Ross \S4.4, ``Necessidade de aporte financeiro e crescimento'', pp.~106--111.}

% =============================================================================
\exsec{15. KMS Designs under 0.25\% share growth --- Berk 18.20--18.22 (18.12--18.15), pp.~560--561}

\asked (18.12) Required production each year and when it first exceeds the 1,100 thousand-unit
capacity. (18.13) Interest and interest tax shields through 2015 when the \$20 million expansion is
delayed to 2012. (18.14) Net income through 2015 with the given depreciation. (18.15) Working capital
needs. (18.20) Free cash flows. (18.21) Continuation value at 8.5 times 2015 EBITDA. (18.22) Firm value
at a 10\% cost of capital.

\approach Rebuild the KMS model (Tables 18.5--18.9, 18.13--18.14) with market share
$10\%+0.25\%\times t$. Value with APV, as in the chapter: unlevered free cash flows and the continuation
value at the 10\% pretax WACC, plus interest tax shields at the 6.8\% debt rate (later tax shields are
inside the multiple). Assumption, stated because the statement says ``adjusted accordingly'': capex is
\$5 million in 2011 (replacement only), \$25 million in 2012 (expansion year) and \$8 million from 2013,
mirroring the chapter's schedule one year later.

\notation
\begin{itemize}
  \item $Q_t=M_t\times s_t$: production (000 units); $S_t=Q_t\times p_t$: sales (\$000); costs $78\%$ of sales; $\tau=35\%$.
  \item $NWC_t$: cash 16\% + receivables 19\% + inventory 20\% $-$ payables 16\% of sales $=39\%$ of sales (2010 actual: 29,207).
  \item $FCF_t=NI_t+(1-\tau)Int_t+Dep_t-\Delta NWC_t-CapEx_t$; $CV_{2015}=8.5\times EBITDA_{2015}$.
\end{itemize}

\step{1}{(18.12) Capacity} Volumes $Q_t=M_t s_t$ (e.g.\ 2012: $11{,}025\times10.5\%=1{,}157.625$):
\begin{center}\small
\begin{tabular}{lrrrrrr}
\toprule
 & 2010 & 2011 & 2012 & 2013 & 2014 & 2015\\
\midrule
Market share (\%) & 10.00 & 10.25 & 10.50 & 10.75 & 11.00 & 11.25\\
Production (000 units) & 1,000.00 & 1,076.25 & 1,157.63 & 1,244.42 & 1,337.05 & 1,435.84\\
\bottomrule
\end{tabular}
\end{center}
\answer{(18.12)} Requirements as in the table; production first exceeds 1,100 in 2012, so the expansion is needed in 2012.

\step{2}{(18.13) Interest and tax shields} By Eq.~(18.1), interest in year $t$ is $6.8\%$ of debt at the
end of $t-1$. Debt is 4,500 through 2011 and $24{,}500$ from the end of 2012:
$0.068(4{,}500)=306$ in 2011 and 2012; $0.068(24{,}500)=1{,}666$ in 2013--2015. Tax shields
$=0.35\times$ interest: $107.1$, $107.1$, $583.1$, $583.1$, $583.1$.
\answer{(18.13)} Interest 306, 306, 1,666, 1,666, 1,666; tax shields 107.1, 107.1, 583.1, 583.1, 583.1 (\$000).

\step{3}{(18.14) Pro forma income statement} Sales $=Q_t p_t$ (2011: $1{,}076.25\times76.51=82{,}343.9$);
then costs, EBITDA, depreciation (given), interest (Step 2), tax 35\%:
\begin{center}\small
\begin{tabular}{lrrrrrr}
\toprule
\$000 & 2010 & 2011 & 2012 & 2013 & 2014 & 2015\\
\midrule
Sales & 74,889.0 & 82,343.9 & 90,341.1 & 99,055.8 & 108,555.1 & 118,916.1\\
Costs except dep.\ (78\%) & $-$58,413.0 & $-$64,228.2 & $-$70,466.0 & $-$77,263.5 & $-$84,673.0 & $-$92,754.5\\
EBITDA & 16,476.0 & 18,115.7 & 19,875.0 & 21,792.3 & 23,882.1 & 26,161.5\\
Depreciation & $-$5,492.0 & $-$5,443.0 & $-$7,398.0 & $-$7,459.0 & $-$7,513.0 & $-$7,561.0\\
EBIT & 10,984.0 & 12,672.7 & 12,477.0 & 14,333.3 & 16,369.1 & 18,600.5\\
Interest & $-$306.0 & $-$306.0 & $-$306.0 & $-$1,666.0 & $-$1,666.0 & $-$1,666.0\\
Pretax income & 10,678.0 & 12,366.7 & 12,171.0 & 12,667.3 & 14,703.1 & 16,934.5\\
Income tax (35\%) & $-$3,737.3 & $-$4,328.3 & $-$4,259.9 & $-$4,433.5 & $-$5,146.1 & $-$5,927.1\\
Net income & 6,940.7 & 8,038.3 & 7,911.2 & 8,233.7 & 9,557.0 & 11,007.4\\
\bottomrule
\end{tabular}
\end{center}
\answer{(18.14)} Net income 8,038.3; 7,911.2; 8,233.7; 9,557.0; 11,007.4 for 2011--2015 (\$000).

\step{4}{(18.15) Working capital} Each item is its percent of sales (2010 actual):
\begin{center}\small
\begin{tabular}{lrrrrrr}
\toprule
\$000 & 2010 & 2011 & 2012 & 2013 & 2014 & 2015\\
\midrule
Cash (16\%) & 11,982.0 & 13,175.0 & 14,454.6 & 15,848.9 & 17,368.8 & 19,026.6\\
Accounts receivable (19\%) & 14,229.0 & 15,645.3 & 17,164.8 & 18,820.6 & 20,625.5 & 22,594.1\\
Inventory (20\%) & 14,978.0 & 16,468.8 & 18,068.2 & 19,811.2 & 21,711.0 & 23,783.2\\
Total current assets & 41,189.0 & 45,289.1 & 49,687.6 & 54,480.7 & 59,705.3 & 65,403.8\\
Accounts payable (16\%) & 11,982.0 & 13,175.0 & 14,454.6 & 15,848.9 & 17,368.8 & 19,026.6\\
Net working capital & 29,207.0 & 32,114.1 & 35,233.0 & 38,631.8 & 42,336.5 & 46,377.3\\
Increase in NWC & & 2,907.1 & 3,118.9 & 3,398.8 & 3,704.7 & 4,040.8\\
\bottomrule
\end{tabular}
\end{center}
\answer{(18.15)} NWC as in the table; increases 2,907.1; 3,118.9; 3,398.8; 3,704.7; 4,040.8 (\$000).

\step{5}{(18.20) Free cash flows} Unlevered net income $=NI+(1-0.35)Int$ ($=0.65\,EBIT$; after-tax
interest is $198.9$ while debt is 4,500 and $1{,}082.9$ after). E.g.\ 2012:
$8{,}110.1+7{,}398-3{,}118.9-25{,}000=-12{,}610.8$.
\begin{center}\small
\begin{tabular}{lrrrrr}
\toprule
\$000 & 2011 & 2012 & 2013 & 2014 & 2015\\
\midrule
Net income & 8,038.3 & 7,911.2 & 8,233.7 & 9,557.0 & 11,007.4\\
Plus: after-tax interest & 198.9 & 198.9 & 1,082.9 & 1,082.9 & 1,082.9\\
Unlevered net income & 8,237.2 & 8,110.1 & 9,316.6 & 10,639.9 & 12,090.3\\
Plus: depreciation & 5,443.0 & 7,398.0 & 7,459.0 & 7,513.0 & 7,561.0\\
Less: increase in NWC & $-$2,907.1 & $-$3,118.9 & $-$3,398.8 & $-$3,704.7 & $-$4,040.8\\
Less: capital expenditures & $-$5,000.0 & $-$25,000.0 & $-$8,000.0 & $-$8,000.0 & $-$8,000.0\\
Free cash flow of firm & 5,773.1 & $-$12,610.8 & 5,376.9 & 6,448.2 & 7,610.6\\
\bottomrule
\end{tabular}
\end{center}
\answer{(18.20)} FCF $=5{,}773.1$; $-12{,}610.8$; $5{,}376.9$; $6{,}448.2$; $7{,}610.6$ (\$000).

\step{6}{(18.21) Continuation value} By Eq.~(18.7),
\begin{align}
  CV_{2015} &= 8.5\times EBITDA_{2015} = 8.5(26{,}161.5) = 222{,}373.0. \label{r2ce15-cv}
\end{align}
\answer{(18.21)} About \$222.4 million at the end of 2015.

\step{7}{(18.22) APV} Discount FCF and $CV$ at 10\% ($1.1^5=1.61051$) and tax shields at 6.8\%:
\begin{align}
  PV(FCF) &= \frac{5{,}773.1}{1.1}+\frac{-12{,}610.8}{1.1^2}+\frac{5{,}376.9}{1.1^3}+\frac{6{,}448.2}{1.1^4}+\frac{7{,}610.6}{1.1^5} \nonumber\\
  &= 5{,}248.3-10{,}422.2+4{,}039.7+4{,}404.2+4{,}725.6 = 7{,}995.6, \nonumber\\
  PV(CV) &= 222{,}373.0/1.61051 = 138{,}076.2, \nonumber\\
  PV(ITS) &= \frac{107.1}{1.068}+\frac{107.1}{1.068^2}+\frac{583.1}{1.068^3}+\frac{583.1}{1.068^4}+\frac{583.1}{1.068^5} \nonumber\\
  &= 100.3+93.9+478.7+448.2+419.6 = 1{,}540.7, \nonumber\\
  V_{2010} &= 7{,}995.6+138{,}076.2+1{,}540.7 = 147{,}612.5. \label{r2ce15-v}
\end{align}
\answer{(18.22)} KMS is worth about \$147.6 million under the 0.25\% scenario, versus \$201.0 million in the
chapter's 1\% scenario (Table 18.14) --- \$53.4 million less (still \$42.5 million less if the chapter's
value is redone at 8.5$\times$: \$190.1 million).

\intuition The value lives in the continuation value (93.5\% of the total), and the continuation value
lives in 2015 EBITDA. Slower share growth cuts 2015 EBITDA from 34,880 to 26,162, and with it most of the
value; the five-year cash flows barely matter. That is the chapter's lesson: growth is valuable only
through the cash flows it adds, and the forecast's horizon multiple must be stress-tested.
\pitfalls
\begin{itemize}
  \item Subtracting the level of NWC instead of its \emph{increase} (Berk's ``Common Mistake'', p.~553).
  \item Discounting the tax shields at 10\%, or adding them twice (net income already contains them; start from \emph{unlevered} income).
  \item Charging the new bond's interest in 2012: under Eq.~(18.1) interest starts the year after the borrowing.
\end{itemize}
\refbib{Berk \S18.3--18.5, pp.~542--556; Tables 18.5--18.9, 18.13--18.14.}

\end{document}
